\documentclass[12pt]{article}
\usepackage[utf8]{inputenc}
\usepackage[T1]{fontenc}
\usepackage{hyperref}
\title{Foau Language Resources for K-12 Curriculum Development}
\author{Compiled by Claude Sonnet 5.5 (claude-sonnet-5-5) for Andr\'as Kornai}
\date{}
\begin{document}
\maketitle

\section{Introduction}
Foau (also Abawiri; ISO 639-3 flh; Glottocode foau1240) is a Lakes Plain language spoken in Fuau village on the Dijai River, a tributary of the Mamberamo, in Papua, Indonesia. Speaker estimates range from about 350 to 400. It is tonal. Very little material exists; this survey is based on a small number of searches.

\section{Primary Educational Resources}
None found. No primers, textbooks or school programmes appeared in the search results.

\section{Linguistic and Reference Materials}
\begin{itemize}
\item Glottolog entry: \url{https://glottolog.org/resource/languoid/id/foau1240}
\item Wikipedia, Abawiri (Foau): \url{https://en.wikipedia.org/wiki/Abawiri}
\item Search summaries mention a 2020 UC Santa Barbara PhD dissertation grammar by Brendon Yoder; no direct URL was seen, so none is given.
\end{itemize}

\section{Religious and Community Resources}
\begin{itemize}
\item Joshua Project profile (Fuau in Indonesia): \url{https://joshuaproject.net/people_groups/11218}
\end{itemize}

\section{Digital and Audio Resources}
\begin{itemize}
\item Endangered Languages Project sample record: \url{https://endangeredlanguages.com/lang/4149/samples/11579}
\end{itemize}
Search summaries note that the Endangered Languages Archive holds recordings with Indonesian and English translations.

\section{Curriculum Development Context}
Foau is spoken by all in the community but is not known to be taught in school. Documentation is mainly academic, and any curriculum would need to be built from the dissertation grammar and archived recordings.
\end{document}